% Part: normal-modal-logic
% Chapter: completeness
% Section: complete-consistent-sets

\documentclass[../../../include/open-logic-section]{subfiles}

\begin{document}

\olfileid{nml}{com}{ccs}

\olsection{完全 $\Sigma$ 一致集}

假设 $\Sigma$ 是一个模态公式的集合——不妨将其视为某个正规模态逻辑的公理或定义性原则。一个集合 $\Gamma$ 是 $\Sigma$ 一致的，当且仅当 $\Gamma \Proves/[\Sigma] \lfalse$，即不存在从 $\Sigma$ 和属于 $\Gamma$ 的各 $!A_i$ 到 $!A_1 \lif (!A_2 \lif \cdots (!A_n \lif
\lfalse)\dots)$ 的推导。我们将构造一个“典范”模型，其中每个世界都被视为一种特殊的 $\Sigma$ 一致集：它不仅仅是 $\Sigma$ 一致的，而且是极大的，即它确定了每个模态公式的真值：对每个 $!A$，要么 $!A \in \Gamma$，要么 $\lnot !A \in \Gamma$：

\begin{defn}
  一个集合 $\Gamma$ 是\emph{完全 $\Sigma$ 一致的}，当且仅当它是 $\Sigma$ 一致的，并且对每个公式~$!A$，要么 $!A \in \Gamma$，要么 $\lnot !A \in \Gamma$。
\end{defn}

完全 $\Sigma$ 一致集 $\Gamma$ 具有许多有用的性质。首先，它在 $\Sigma$ 中是演绎封闭的，即如果 $\Gamma \Proves[\Sigma] !A$，那么 $!A \in \Gamma$。这尤其意味着 $\Sigma$ 中每个公式的所有实例也都属于 $\Gamma$。此外，$\Gamma$ 中的成员关系反映了命题联结词的真值条件。当我们定义“典范模型”时，这一点将非常重要。

\begin{prop}\ollabel{prop:ccs-properties}
  假设 $\Gamma$ 是完全 $\Sigma$ 一致的，那么：
  \begin{enumerate}
  \item \ollabel{prop:ccs-closed}%
    $\Gamma$ 在 $\Sigma$ 中是演绎封闭的。
  \item \ollabel{prop:ccs-sigma}%
    $\Sigma \subseteq \Gamma$。
  \tagitem{prvFalse}{\ollabel{prop:ccs-lfalse}%
    $\lfalse \notin \Gamma$}{}
  \tagitem{prvTrue}{\ollabel{prop:ccs-ltrue}%
    $\ltrue \in \Gamma$}{}
  \item \ollabel{prop:ccs-lnot}%
    $\lnot!A \in \Gamma$ 当且仅当 $!A \notin \Gamma$。
  \tagitem{prvAnd}{\ollabel{prop:ccs-land}%
    $!A \land !B \in \Gamma$ 当且仅当 $!A \in \Gamma$ 且 $!B \in \Gamma$}{}
  \tagitem{prvOr}{\ollabel{prop:ccs-lor}%
    $!A \lor !B \in \Gamma$ 当且仅当 $!A \in \Gamma$ 或 $!B \in \Gamma$}{}
  \tagitem{prvIf}{\ollabel{prop:ccs-lif}%
    $!A \lif !B \in \Gamma$ 当且仅当 $!A \notin \Gamma$ 或 $!B \in \Gamma$}{}
  \tagitem{prvIff}{\ollabel{prop:ccs-liff}%
    $!A \liff !B \in \Gamma$ 当且仅当要么 $!A \in \Gamma$ 且 $!B \in \Gamma$，要么 $!A \notin \Gamma$ 且 $!B \notin \Gamma$}{}
  \end{enumerate}
\end{prop}

\begin{proof}
  \begin{enumerate}
  \item 假设 $\Gamma \Proves[\Sigma] !A$ 但 $!A \notin
    \Gamma$。那么由于 $\Gamma$ 是完全 $\Sigma$ 一致的，
    $\lnot!A \in \Gamma$。这将导致 $\Gamma$ 不一致，因为
    $!A, \lnot !A \Proves[\Sigma] \lfalse$。

  \item 如果 $!A \in \Sigma$ 那么 $\Gamma \Proves[\Sigma] !A$，并且由演绎封闭性可知 $!A
    \in \Gamma$，即由第~\olref{prop:ccs-closed}~项。

  \tagitem{prvFalse}{若 $\lfalse \in \Gamma$，则 $\Gamma
    \Proves[\Sigma] \lfalse$，因此 $\Gamma$ 将是
    $\Sigma$-不一致的。}{}

  \tagitem{prvTrue}{$\Gamma \Proves[\Sigma] \ltrue$，故由演绎封闭性得 $\ltrue \in
    \Gamma$，即由第~\olref{prop:ccs-closed}~项。}{}

\item 若 $\lnot!A \in \Gamma$，则由一致性可知 $!A \notin
    \Gamma$；且若 $!A \notin \Gamma$，则由于 $\Gamma$ 是完全 $\Sigma$ 一致的，有 $\lnot!A \in \Gamma$。

  \tagitem{prvAnd}{\iftag{probAnd}{练习。}{假设 $!A \land !B \in
      \Gamma$。由于 $(!A \land !B) \lif !A$ 是一个重言式实例，故由演绎封闭性（即
      第~\olref{prop:ccs-closed}~项）得 $!A \in \Gamma$。类似可得 $!B \in \Gamma$。
      另一方面，假设 $!A \in \Gamma$ 且 $!B \in
      \Gamma$。那么由演绎封闭性可得 $(!A \land !B) \in
      \Gamma$，因为 $!A \lif (!B \lif (!A \land !B))$ 是一个
      重言式实例。}}{}

  \tagitem{prvOr}{\iftag{probOr}{练习。}{假设 $!A \lor !B \in
      \Gamma$，且 $!A \notin \Gamma$ 以及 $!B \notin \Gamma$。由于
      $\Gamma$ 是完全 $\Sigma$ 一致的，故 $\lnot!A \in \Gamma$
      且 $\lnot !B \in \Gamma$。那么 $\lnot(!A \lor !B) \in \Gamma$，
      因为 $\lnot !A \lif (\lnot !B \lif \lnot (!A \lor !B))$ 是一个
      重言式实例。这将意味着 $\Gamma$ 是
      $\Sigma$-不一致的，矛盾。}}{}

  \tagitem{prvIf}{\iftag{probIf}{练习。}{假设 $!A \lif !B \in
      \Gamma$ 且 $!A \in \Gamma$；则 $\Gamma \Proves[\Sigma] !B$，
      从而由演绎封闭性得 $!B \in \Gamma$。反之，若 $!A
      \lif !B \notin \Gamma$，则由于 $\Gamma$ 是完全
      $\Sigma$ 一致的，有 $\lnot (!A \lif !B) \in \Gamma$。由于
      $\lnot(!A \lif !B) \lif !A$ 是重言式实例，故由演绎封闭性得 $!A \in
      \Gamma$。由于 $\lnot(!A \lif !B) \lif
      \lnot !B$ 是重言式实例，故得 $\lnot !B \in
      \Gamma$。那么由于 $\Gamma$ 是
      $\Sigma$-一致的，有 $!B \notin \Gamma$。}}{}

  \tagitem{prvIff}{\iftag{probIff}{练习。}{假设 $!A \liff !B \in
      \Gamma$。若 $!A \in \Gamma$，则 $!B \in \Gamma$，因为 $(!A
      \liff !B) \lif (!A \lif !B)$ 是重言式
      实例。类似地，若 $!B \in \Gamma$，则 $!A \in \Gamma$。
      因此，要么 $!A \in \Gamma$ 且 $!B \in \Gamma$，要么 $!A \notin \Gamma$ 且 $!B \notin \Gamma$。

      反之，假设 $!A \liff !B \notin \Gamma$。由于 $\Gamma$
      是完全 $\Sigma$ 一致的，故 $\lnot (!A \liff !B) \in
      \Gamma$。由于 $\lnot(!A \liff !B) \lif (!A \lif \lnot !B)$ 是一个
      重言式实例，如果 $!A \in \Gamma$ 则 $\lnot !B \in
      \Gamma$，并且由于 $\Gamma$ 是 $\Sigma$-一致的，故 $!B \notin
      \Gamma$。类似地，如果 $!B \in \Gamma$ 则 $!A \notin \Gamma$。
      因此，不可能出现 $!A \in \Gamma$ 且 $!B \in \Gamma$，也不可能出现 $!A \notin
      \Gamma$ 且 $!B \notin \Gamma$。}}{}
  \end{enumerate}
\end{proof}

\begin{probtag}{probAnd,probOr,probIf,probIff}
  请完成 \olref[nml][com][ccs]{prop:ccs-properties} 的证明。
\end{probtag}

\end{document}